\chapter{The Ledger of Corrections}

The raw recoil is not yet \(\alpha\), and it is not even yet \(h/m_X\). It is a smeared physical reading, shaped
by every gauge and every nuisance field the apparatus carried. The user must now do what the construction has
been teaching from the first difference: tange what must be separated, funge what can be pooled, and refuse to
let a correction become invisible merely because it is inconvenient.

\section{Tange the Biases}

A bias is a movable label that has learned to push the count. It may be a laser wavefront whose curvature makes
one part of the atom cloud see a different kick from another. It may be an AC Stark shift from light intensity,
a Zeeman shift from magnetic field, a Coriolis term from Earth's rotation, a gravity-gradient term from the local
mass distribution, a mean-field shift from atom density, or a line-shape bias from the fitting model. Each bias
must be tanged first: selected by name, given its own symbol, its own sign, its own method of estimation.

The correction ledger should have one row per bias:
\[
\begin{array}{c|c|c|c}
\hbox{effect} & \hbox{correction to } h/m_X & \hbox{standard uncertainty} & \hbox{evidence}\\
\end{array}
\]
The evidence column matters. A correction obtained by auxiliary measurement is not the same as one obtained by
simulation, and neither is the same as one bounded by deliberately exaggerating the effect. The apparatus can use
all three, but it must label them. A simulated wavefront correction may be excellent; it is still a simulated
correction. A density extrapolation may be clean; it is still an extrapolation. The grade travels with the row.

Every sign convention should be checked by reversal. Reverse the recoil direction. Reverse the magnetic field
where possible. Reverse the launch direction. Swap pulse order if the interferometer allows it. A real recoil
term changes in the way the geometry says it should; a hidden bias often does not. Reversal is the apparatus'
oldest test in a laboratory coat: relabel what the hand is free to relabel and watch what survives. What
survives belongs to the world. What follows the relabeling belongs to the hand, and must be accounted for as
gauge.

\section{Correction Rows in Lean}

In the companion file, each correction row is a record. The name says what moved the recoil, \(\texttt{delta}\)
is the signed correction to \(h/m_X\), \(\texttt{sigma}\) is the standard uncertainty assigned to that correction,
and \(\texttt{evidence}\) records how the row was earned:

\begin{quote}\small
\begin{verbatim}
structure Correction where
  name : String
  delta : Float
  sigma : Float
  evidence : String
\end{verbatim}
\end{quote}

An example correction list looks like this:

\begin{quote}\small
\begin{verbatim}
def corrections : List Correction :=
  [ { name := "wavefront"
      delta := 0.000000003e-9
      sigma := 0.000000002e-9
      evidence := "auxiliary wavefront map" },
    { name := "AC Stark"
      delta := -0.000000001e-9
      sigma := 0.000000001e-9
      evidence := "power reversal" } ]
\end{verbatim}
\end{quote}

The list is deliberately plain. No row hides behind a paragraph. No correction enters without a sign. No
uncertainty enters without a value. No evidence enters without a sentence the later reader can challenge. The
printed table and the Lean list should agree row for row.

\section{Funge the Repetitions}

Once the biases are named, the repeated runs can be pooled. Pool only like with like. Runs taken under the same
carrier preparation, same gauge, same correction model, and same blind offset may be funged into a mean. Runs
taken under changed conditions should first be held apart and compared. The question is not merely whether all
the digits can be averaged. The question is whether the residuals behave as noise after the tanged biases have
been removed.

For each block of runs compute the corrected recoil ratio
\[
  \left(\frac{h}{m_X}\right)_{\rm corr}
  =
  \left(\frac{h}{m_X}\right)_{\rm raw}
  + \sum_i \Delta_i,
\]
where \(\Delta_i\) are the signed corrections. Then inspect the residuals:
\[
  r_j =
  \left(\frac{h}{m_X}\right)_{j,{\rm corr}}
  -
  \left\langle\frac{h}{m_X}\right\rangle .
\]
If the residuals are wider than their stated noise, do not tighten the error bar by force. The apparatus has
found a leftover. Either identify the missing tanged effect or enlarge the uncertainty honestly. A beautiful mean
with a dishonest scatter is not a measurement; it is an argument with the data concealed.

The pooling should be hierarchical. First pool shots into fringes. Then pool fringes into runs. Then pool runs
into blocks. Then pool blocks only after the block-to-block consistency is tested. At every level, the apparatus
asks the same question: does repetition return the same answer within the noise the tape says it should? If yes,
funge. If no, tange the difference and find out what moved.

Lean represents the last pooling step as a correction closure:

\begin{quote}\small
\begin{verbatim}
def correctionSum : List Correction -> Float
  | [] => 0.0
  | c :: cs => c.delta + correctionSum cs

def correctionVariance : List Correction -> Float
  | [] => 0.0
  | c :: cs => square c.sigma + correctionVariance cs

def closeRecoil (raw : Datum)
    (corrections : List Correction) : ClosedRecoil :=
  { raw := raw
    corrections := corrections
    value := raw.value + correctionSum corrections
    sigma := Float.sqrt
      (square raw.sigma + correctionVariance corrections) }
\end{verbatim}
\end{quote}

This is a deliberately modest model: it treats the correction uncertainties as independent. When correlations
matter, replace the variance rule with the covariance-aware rule and let the compiler show every place the change
must travel. The essential point remains the same: the corrected recoil is not hand-edited. It is constructed
from the raw recoil and the frozen correction list.

The file includes a tiny proof of that construction:

\begin{quote}\small
\begin{verbatim}
theorem closeRecoil_value
    (raw : Datum) (corrections : List Correction) :
    (closeRecoil raw corrections).value =
      raw.value + correctionSum corrections := by
  rfl
\end{verbatim}
\end{quote}

The proof is small because the definition is honest. If someone later changes \(\texttt{closeRecoil}\) so that
the value is no longer raw plus corrections, this theorem is where the lie fails to typecheck.

\section{Carry the Uncertainty}

Uncertainty is not embarrassment. It is the residue carried honestly. The final \(h/m_X\) needs a combined
standard uncertainty with statistical and systematic parts shown separately before they are combined. Statistical
uncertainty comes from repeated readings. Systematic uncertainty comes from correction rows, calibration rows,
and model choices. Correlations must be named. Two corrections that share a laser calibration do not become
independent because they sit on different lines of the table.

The propagation into \(\alpha\) is simple in form. Since
\[
  \alpha =
  \left(
  \frac{2R_\infty}{c}\,
  \frac{A_{\rm r}(X)}{A_{\rm r}(e)}\,
  \frac{h}{m_X}
  \right)^{1/2},
\]
the fractional uncertainty, ignoring correlations among the imported constants, is
\[
\begin{aligned}
  \left(\frac{2u(\alpha)}{\alpha}\right)^2
  &=
    \left(\frac{u(h/m_X)}{h/m_X}\right)^2
    +
    \left(\frac{u(R_\infty)}{R_\infty}\right)^2 \\
  &\quad+
    \left(\frac{u(A_{\rm r}(X))}{A_{\rm r}(X)}\right)^2
    +
    \left(\frac{u(A_{\rm r}(e))}{A_{\rm r}(e)}\right)^2 .
\end{aligned}
\]
The exact analysis should include the correlations supplied by the constants adjustment if they matter at the
target precision. The apparatus' rule is simple: correlations are not optional once they are large enough to
move the last quoted digit.

In Lean, the corresponding independent-input rule is:

\begin{quote}\small
\begin{verbatim}
def alphaRelativeSigma (input : AlphaInput) : Float :=
  let e := input.constants.relativeMassElectron
  0.5 * Float.sqrt
    (square (recoilRelativeSigma input.recoil) +
     square (relativeSigma input.constants.rydberg) +
     square (relativeSigma input.constants.relativeMassX) +
     square (relativeSigma e))
\end{verbatim}
\end{quote}

The factor \(0.5\) is the square-root bridge. The product computes \(\alpha^2\), so the fractional uncertainty
of \(\alpha\) is half the fractional uncertainty of the product. Again, this is not a license to ignore
correlation. It is the base case, the clean surface from which the more careful covariance calculation can be
introduced.

\section{Close the Ledger Before Opening the Blind}

The correction table, pooling rule, rejection criteria, uncertainty propagation, and comparison tests should be
frozen before the blind offset is removed. This is the human side of the gauge. If the operator can still change
which runs count after seeing whether the answer agrees with CODATA, the operator has become an unrecorded
correction. The apparatus cannot prevent that by logic. It can only require the tape: freeze the ledger, open the
blind, compute once, and leave the record of the opening where a reader can see it.

When the blind is opened, the corrected \(h/m_X\) becomes the measured input to the alpha bridge. Nothing in the
systematic table should change merely because the resulting \(\alpha\) is surprising. A surprising answer may
send the user back to run the experiment again, or to build a better correction, or to compare against an
independent method. It must not send the user back to quietly edit the past. The past is the tape. The tape is
the measurement's memory.
